% Generated by code/assemble.py; edit part_B_audit.tex.
\subsection{The unique predecessor and scalar tests}

\begin{result}{B.101}{Proved}{Two exceptional-prime tests select a unique class-two class}
Let $(A,B)\in\mathcal G_{\rm ord}$ satisfy
\[
(v_3(A),v_3(B))=(1,0),\qquad
(v_2(A),v_2(B))=(0,1).
\]
Use the actual integral class-two arithmetic quotient of the tripod,
the tangential splitting at zero, and geometric conjugacy of continuous
sections. There is exactly one class-two section class above $(A,B)$
whose third obstruction vanishes at both of the following pairs of
rational place and coefficient prime:
\[
(q,\ell)=(3,2),\qquad(q,\ell)=(2,3).
\]
Consequently at most one class-two section class above this pair can
be the projection of an integral class-three section. In particular
the twenty-four alternatives in the original A.44 family reduce to
one canonical candidate before the other local or continuous tests
are imposed. The assertion does not say that this candidate lifts.
\end{result}
\begin{proof}
We use the independently checked weight-two calculation
$H^1(G_{\mathbb Q},\widehat{\mathbb Z}(2))=\mathbb Z/24$, generated
by the class of $\eta(g)=(\chi(g)^2-1)/24$, together with the actual
class-two lifting result for ordinary pairs. For fixed literal abelian
cocycles $a,b$, write a section in Hall coordinates as
$y^{b(g)}x^{a(g)}C^{h(g)}s_0(g)$. Its equation is
$\delta h=-a\cup b$. Every other class-two class is represented by
$h+k\eta$, for one $k\in\mathbb Z/24$. Central coboundaries account
for the change of representative. There is no additional abelian
stabilizer: $H^0(G_{\mathbb Q},\widehat{\mathbb Z}(1)^2)=0$,
since each local factor is torsion-free and a nonidentity cyclotomic
scalar annihilates no nonzero vector.

At $q=3$, with coefficient prime two, ordinarity and $v_3(A)=1$
make the local class of $B$ that of the principal unit $1-x_3$,
and hence zero in $H^1(G_{\mathbb Q_3},\mathbb Z_2(1))$.
Thus the literal cocycle $b$ is $\delta u$ for a unique
$u\in\mathbb Z_2(1)$. For tame inertia $T$ and arithmetic Frobenius
$F$, put $\alpha=a(F)$ and $\gamma=h(F)$. We have
\[
a(T)=1,\quad b(T)=0,\quad b(F)=2u,\quad h(T)=u.
\]
The last equality follows from the actual class-two tame relation;
the scalar $q-1$ is not a zero divisor in the integral module, so
there is no extra inertia freedom. Substitution into the actual third
residue formula gives
\[
Z_D=\gamma-3\alpha u,\qquad
Z_F=-8\binom{u+1}{2}.
\]
Here we use the actual cochain and purity comparison of B.22--B.23.
The local obstruction group is
$H^2(G_{\mathbb Q_3},\mathbb Z_2(3))=\mathbb Z_2/8$ in each
coordinate. The second displayed residue is already zero there.
Twisting $h$ by $k\eta$ leaves its inertia unchanged and changes
$\gamma$ by $k\eta(F)=k/3$. Multiplication by $1/3$ is a unit
modulo eight. Exactly one $k\pmod8$ therefore makes both residues
zero.

At $q=2$, with coefficient prime three, the ordinary local class of
$A=1-(1-x_2)$ is a principal two-adic unit and is zero in the
three-primary Kummer group. Write $a=\delta v$, uniquely, and set
$\beta=b(F)$ and $\gamma=h(F)$ in this component. The tame data are
\[
a(T)=0,\quad b(T)=1,\quad a(F)=v,\quad h(T)=-v.
\]
The residue formulas now give
\[
Z_D=3\binom{v+1}{2},\qquad
Z_F=\gamma+\beta v-\frac18.
\]
The local obstruction group in each coordinate is $\mathbb Z_3/3$.
Again the first residue already vanishes, and twisting changes the
second by $k\eta(F)=k/8$, a unit multiple of $k$ modulo three.
Exactly one $k\pmod3$ makes this local obstruction zero.

The two congruences select exactly one $k\pmod{24}$ by the Chinese
remainder theorem. All formulas are insensitive to central coboundaries,
which change the Frobenius values by multiples of $q^2-1$.
Changing aligned cocycle representatives amounts to geometric
conjugacy and preserves the obstruction. This proves existence and
uniqueness for these two tests, and the necessity for a global
class-three lift follows by continuous restriction to the two local
groups. The polynomial substitutions and the resulting bijection
$k\mapsto(k/3\bmod8,k/8\bmod3)$ are checked in
\texttt{code/ghost\_audit\_checks.py}.
\end{proof}


\begin{result}{B.102}{Proved}{The remaining fresh-support test in two scalar sequences}
For an original A.44 pair let $s$ be the canonical class-two candidate
of B.101 and put $n_j=(j+4)!$. Fix a coefficient prime $\ell$.
At each positive $A$-support prime $q_j\ne\ell$, write
$b=\delta u_j$ locally, and use Frobenius coordinates
$\alpha_j=a(F_{q_j})$, $\gamma_j=h(F_{q_j})$.
At each positive $B$-support prime $r_j\ne\ell$, write
$a=\delta v_j$ and put
$\beta_j=b(F_{r_j})$, $\widetilde\gamma_j=h(F_{r_j})$.
The two potentially nonintegral normalized residues are
\[
T_{A,j}=\frac{n_j(\gamma_j-q_j\alpha_j u_j)}{q_j^2-1},
\qquad
T_{B,j}=\frac{n_j(\widetilde\gamma_j+\beta_jv_j)}{r_j^2-1}.
\]
The fresh-support local class-three obstructions all vanish if and
only if these numbers belong to $\mathbb Z_\ell$ for every indicated
index and every $\ell$. After the finitely many old-place residues
have also been tested, the global continuous class-three obstruction
vanishes if and only if, in addition,
$T_{A,j}\to0$ and $T_{B,j}\to0$ in $\mathbb Z_\ell$ for every $\ell$.
The quotient formulas are exact, including coefficient primes two
and three. This is a reduction of B1 to explicit arithmetic conditions
on one section class, not a construction of a prime sequence satisfying
them.
\end{result}
\begin{proof}
We use the strict continuous localization calculation of A.22 and
the literal cochain comparison B.22--B.23, in
which the third obstruction is a class in
$R_\ell/(q^2-1)_qR_\ell$ for each of its two central coordinates,
where $R_\ell$ consists of families tending to zero. Local vanishing
is coordinatewise divisibility by $q^2-1$; global vanishing requires
the divided family itself to lie in $R_\ell$.

At $q=q_j$ the data are
$m=n_j$, $n=0$, $t=n_ju_j$, and $\beta=(q-1)u_j$.
Direct substitution into the two third-residue polynomials gives
\[
\frac{Z_{D,q}}{q^2-1}
 =T_{A,j}+u_j\binom{n_j}{2},\qquad
\frac{Z_{F,q}}{q^2-1}
 =-n_j\binom{u_j+1}{2}.
\]
At $q=r_j$ we instead have $m=0$, $n=n_j$, $t=-n_jv_j$,
and $\alpha=(q-1)v_j$. The exact identities are
\[
\frac{Z_{D,q}}{q^2-1}
 =n_j\binom{v_j+1}{2},\qquad
\frac{Z_{F,q}}{q^2-1}
 =T_{B,j}-\frac{n_j}{24}-v_j\binom{n_j}{2}.
\]
The term $n_j/24$ retains the actual Ihara-action correction; it
must not be discarded at two or three. It is integral because every
original factorial $n_j$ is divisible by twenty-four.

Each $u_j,v_j$ belongs to the integral local Tate module after the
chosen trivialization. The binomial expression $z(z+1)/2$ is integral
on every $\mathbb Z_\ell$, including $\mathbb Z_2$. Since
$n_j\to0$ in every $\mathbb Z_\ell$, all the other terms in the
four displayed quotients tend to zero: the only possible fixed
denominators are two and twenty-four. They are therefore irrelevant
both to divisibility and, after integrality is established, to the
zero-limit condition. This proves the equivalence at fresh support.
At zero-valuation primes the integral class-two tame relation forces
all inertia coordinates to vanish, so both residues are zero. The
old primes two and three contribute finitely many terms, which must
be tested for integrality but do not alter convergence of the family.
A.22 then gives the asserted global equivalence. No local section
theorem enters any of these computations.
\end{proof}

